
\documentclass[a4paper,11pt]{article}
\usepackage[a4paper, margin=8em]{geometry}

% usa i pacchetti per la scrittura in italiano
\usepackage[french,italian]{babel}
\usepackage[T1]{fontenc}
\usepackage[utf8]{inputenc}
\frenchspacing 

% usa i pacchetti per la formattazione matematica
\usepackage{amsmath, amssymb, amsthm, amsfonts}

% usa altri pacchetti
\usepackage{gensymb}
\usepackage{hyperref}
\usepackage{standalone}

% imposta il titolo
\title{Appunti Elettronica Digitale}
\author{Luca Seggiani}
\date{2026}

% imposta lo stile
% usa helvetica
\usepackage[scaled]{helvet}
% usa palatino
\usepackage{palatino}
% usa un font monospazio guardabile
\usepackage{lmodern}

\renewcommand{\rmdefault}{ppl}
\renewcommand{\sfdefault}{phv}
\renewcommand{\ttdefault}{lmtt}

% disponi il titolo
\makeatletter
\renewcommand{\maketitle} {
	\begin{center} 
		\begin{minipage}[t]{.8\textwidth}
			\textsf{\huge\bfseries \@title} 
		\end{minipage}%
		\begin{minipage}[t]{.2\textwidth}
			\raggedleft \vspace{-1.65em}
			\textsf{\small \@author} \vfill
			\textsf{\small \@date}
		\end{minipage}
		\par
	\end{center}

	\thispagestyle{empty}
	\pagestyle{fancy}
}
\makeatother

% disponi teoremi
\usepackage{tcolorbox}
\newtcolorbox[auto counter, number within=section]{theorem}[2][]{%
	colback=blue!10, 
	colframe=blue!40!black, 
	sharp corners=northwest,
	fonttitle=\sffamily\bfseries, 
	title=~\thetcbcounter: #2, 
	#1
}

% disponi definizioni
\newtcolorbox[auto counter, number within=section]{definition}[2][]{%
	colback=red!10,
	colframe=red!40!black,
	sharp corners=northwest,
	fonttitle=\sffamily\bfseries,
	title=~\thetcbcounter: #2,
	#1
}

% U.D.M
\newcommand{\amp}{\ensuremath{\mathrm{A}}}
\newcommand{\volt}{\ensuremath{\mathrm{V}}}
\newcommand{\meter}{\ensuremath{\mathrm{m}}}
\newcommand{\second}{\ensuremath{\mathrm{s}}}
\newcommand{\farad}{\ensuremath{\mathrm{F}}}
\newcommand{\henry}{\ensuremath{\mathrm{H}}}
\newcommand{\siemens}{\ensuremath{\mathrm{S}}}

% circuiti
\usepackage{circuitikz}
\usetikzlibrary{babel}

% disegni
\usepackage{pgfplots}
\pgfplotsset{width=10cm,compat=1.9}

% disponi codice
\usepackage{listings}
\usepackage[table]{xcolor}

\definecolor{codegreen}{rgb}{0,0.6,0}
\definecolor{codegray}{rgb}{0.5,0.5,0.5}
\definecolor{codepurple}{rgb}{0.58,0,0.82}
\definecolor{backcolour}{rgb}{0.95,0.95,0.92}

\lstdefinestyle{codestyle}{
		backgroundcolor=\color{black!5}, 
		commentstyle=\color{codegreen},
		keywordstyle=\bfseries\color{magenta},
		numberstyle=\sffamily\tiny\color{black!60},
		stringstyle=\color{green!50!black},
		basicstyle=\ttfamily\footnotesize,
		breakatwhitespace=false,         
		breaklines=true,                 
		captionpos=b,                    
		keepspaces=true,                 
		numbers=left,                    
		numbersep=5pt,                  
		showspaces=false,                
		showstringspaces=false,
		showtabs=false,                  
		tabsize=2
}

\lstdefinestyle{shellstyle}{
		backgroundcolor=\color{black!5}, 
		basicstyle=\ttfamily\footnotesize\color{black}, 
		commentstyle=\color{black}, 
		keywordstyle=\color{black},
		numberstyle=\color{black!5},
		stringstyle=\color{black}, 
		showspaces=false,
		showstringspaces=false, 
		showtabs=false, 
		tabsize=2, 
		numbers=none, 
		breaklines=true
}

\lstdefinelanguage{javascript}{
	keywords={typeof, new, true, false, catch, function, return, null, catch, switch, var, if, in, while, do, else, case, break},
	keywordstyle=\color{blue}\bfseries,
	ndkeywords={class, export, boolean, throw, implements, import, this},
	ndkeywordstyle=\color{darkgray}\bfseries,
	identifierstyle=\color{black},
	sensitive=false,
	comment=[l]{//},
	morecomment=[s]{/*}{*/},
	commentstyle=\color{purple}\ttfamily,
	stringstyle=\color{red}\ttfamily,
	morestring=[b]',
	morestring=[b]"
}

% disponi sezioni
\usepackage{titlesec}

\titleformat{\section}
	{\sffamily\Large\bfseries} 
	{\thesection}{1em}{} 
\titleformat{\subsection}
	{\sffamily\large\bfseries}   
	{\thesubsection}{1em}{} 
\titleformat{\subsubsection}
	{\sffamily\normalsize\bfseries} 
	{\thesubsubsection}{1em}{}

% disponi alberi
\usepackage{forest}

\forestset{
	rectstyle/.style={
		for tree={rectangle,draw,font=\large\sffamily}
	},
	roundstyle/.style={
		for tree={circle,draw,font=\large}
	}
}

% disponi algoritmi
\usepackage{algorithm}
\usepackage{algorithmic}
\makeatletter
\renewcommand{\ALG@name}{Algoritmo}
\makeatother

% disponi numeri di pagina
\usepackage{fancyhdr}
\fancyhf{} 
\fancyfoot[L]{\sffamily{\thepage}}

\makeatletter
\fancyhead[L]{\raisebox{1ex}[0pt][0pt]{\sffamily{\@title \ \@date}}} 
\fancyhead[R]{\raisebox{1ex}[0pt][0pt]{\sffamily{\@author}}}
\makeatother

\begin{document}
% sezione (data)
\section{Lezione del 11-05-26}

% stili pagina
\thispagestyle{empty}
\pagestyle{fancy}

% testo
\subsubsection{Sintesi della porta XOR in CMOS}
Vediamo, a scopo di esempio, la sintesi in tecnologia CMOS della porta XOR a 2 ingressi. Prendiamo l'equazione dello XOR come:
$$
y = A \overline{B} + \overline{A} B
$$
Ricordiamo che le forme che volevamo ottenere erano:
$$
y = f \left( \overline{A}, \overline{B}, ... \right) \quad \text{(PUN)}
$$
$$
\overline{y} = f \left( A, B, ... \right) \quad \text{(PDN)}
$$

L'equazione della $y$ dello XOR data non è in nessuna delle due forme, per cui prendiamo la negazione:
$$
\overline{y} = \overline{ A \overline{B} + \overline{A} B } = \overline{ A \overline{B} } \cdot \overline{ \overline{A} B } = \left( \overline{A} + B \right) \cdot \left( A + \overline{B} \right) = \overline{A} A + \overline{A} \overline{B} + B A + B \overline{B} = \overline{A} \overline{B} + B A
$$
che ancora non è nella forma richiesta. 

Quello che saremo costretti a fare sarà allora prendere il circuito:
\begin{center}
\begin{tikzpicture}[transform shape]
	% Paths, nodes and wires:
	\node[pmos] at (3, 5.77){};
	\node[pmos, xscale=-1, yscale=-1] at (6, 5.77){};
	\node[pmos] at (3, 3.77){};
	\node[pmos, xscale=-1, yscale=-1] at (6, 3.77){};
	\draw (3, 5) -| (3, 4.54);
	\draw (6, 5) -- (6, 4.54);
	\draw (3, 6.54) -- (6, 6.54);
	\draw (3, 3) -- (6, 3);
	\node[nmos, xscale=-1, yscale=-1] at (6.02, 1.25){};
	\node[nmos] at (3.02, 1.25){};
	\node[nmos, xscale=-1, yscale=-1] at (6.02, -0.75){};
	\node[nmos] at (3.02, -0.75){};
	\draw (3.02, 0.48) -| (3.02, 0.02);
	\draw (6.02, 0.48) -- (6.02, 0.02);
	\draw (6.02, 2.02) |- (3.02, 2.02);
	\draw (4.5, 6.54) -| (4.5, 6.79);
	\node[vcc] at (4.5, 6.79){};
	\node[ground] at (3.02, -1.52){};
	\node[ground] at (6.02, -1.52){};
	\draw (4.48, 3) -| (4.5, 2.02);
	\draw[-stealth] (4.5, 2.5) -- (7, 2.5);
	\node[shape=rectangle, minimum width=0.465cm, minimum height=0.465cm] at (7.25, 2.5){} node[anchor=west, align=left, text width=0.077cm, inner sep=6pt] at (7, 2.5){$y$};
	\node[shape=rectangle, minimum width=0.465cm, minimum height=0.465cm] at (2, 6.25){} node[anchor=center, align=center, text width=0.077cm, inner sep=6pt] at (2, 6.25){$\overline{A}$};
	\node[shape=rectangle, minimum width=0.465cm, minimum height=0.465cm] at (7, 6.25){} node[anchor=center, align=center, text width=0.077cm, inner sep=6pt] at (7, 6.25){$A$};
	\node[shape=rectangle, minimum width=0.465cm, minimum height=0.465cm] at (2, 4.25){} node[anchor=center, align=center, text width=0.077cm, inner sep=6pt] at (2, 4.25){$B$};
	\node[shape=rectangle, minimum width=0.465cm, minimum height=0.465cm] at (7, 4.25){} node[anchor=center, align=center, text width=0.077cm, inner sep=6pt] at (7, 4.25){$\overline{B}$};
	\node[shape=rectangle, minimum width=0.465cm, minimum height=0.465cm] at (2, 1.75){} node[anchor=center, align=center, text width=0.077cm, inner sep=6pt] at (2, 1.75){$\overline{A}$};
	\node[shape=rectangle, minimum width=0.465cm, minimum height=0.465cm] at (7, 1.75){} node[anchor=center, align=center, text width=0.077cm, inner sep=6pt] at (7, 1.75){$A$};
	\node[shape=rectangle, minimum width=0.465cm, minimum height=0.465cm] at (2, -0.25){} node[anchor=center, align=center, text width=0.077cm, inner sep=6pt] at (2, -0.25){$\overline{B}$};
	\node[shape=rectangle, minimum width=0.465cm, minimum height=0.465cm] at (7, -0.25){} node[anchor=center, align=center, text width=0.077cm, inner sep=6pt] at (7, -0.25){$B$};
\end{tikzpicture}
\end{center}
cioè prevedere una rete con entrate non negate (alla PUN), e negate (alla PDN). L'unico modo che avremo per creare questa rete è di inserire una coppia di inverter, che fornisca la negazione sia di $A$ che di $B$.
Questo ha pressappoco il seguente aspetto:
\begin{center}
\begin{tikzpicture}[transform shape]
	% Paths, nodes and wires:
	\node[pmos] at (3, 5.77){};
	\node[pmos, xscale=-1, yscale=-1] at (6, 5.77){};
	\node[pmos] at (3, 3.77){};
	\node[pmos, xscale=-1, yscale=-1] at (6, 3.77){};
	\draw (3, 5) -| (3, 4.54);
	\draw (6, 5) -- (6, 4.54);
	\draw (3, 6.54) -- (6, 6.54);
	\draw (3, 3) -- (6, 3);
	\node[nmos, xscale=-1, yscale=-1] at (6, 1.25){};
	\node[nmos] at (3, 1.25){};
	\node[nmos, xscale=-1, yscale=-1] at (6, -0.75){};
	\node[nmos] at (3, -0.75){};
	\draw (3, 0.48) -| (3, 0.02);
	\draw (6, 0.48) -- (6, 0.02);
	\draw (6, 2.02) |- (3, 2.02);
	\draw (4.5, 6.54) -| (4.5, 6.79);
	\node[vcc] at (4.5, 6.79){};
	\node[ground] at (3, -1.52){};
	\node[ground] at (6, -1.52){};
	\draw (4.48, 3) -| (4.5, 2.02);
	\draw[-stealth] (4.5, 2.5) -- (7, 2.5);
	\node[shape=rectangle, minimum width=0.465cm, minimum height=0.465cm] at (7.25, 2.5){} node[anchor=west, align=left, text width=0.077cm, inner sep=6pt] at (7, 2.5){$y$};
	\node[shape=rectangle, minimum width=0.465cm, minimum height=0.465cm] at (2, 6.25){} node[anchor=center, align=center, text width=0.077cm, inner sep=6pt] at (2, 6.25){$\overline{A}$};
	\node[shape=rectangle, minimum width=0.465cm, minimum height=0.465cm] at (7, 6.25){} node[anchor=center, align=center, text width=0.077cm, inner sep=6pt] at (7, 6.25){$A$};
	\node[shape=rectangle, minimum width=0.465cm, minimum height=0.465cm] at (2, 4.25){} node[anchor=center, align=center, text width=0.077cm, inner sep=6pt] at (2, 4.25){$B$};
	\node[shape=rectangle, minimum width=0.465cm, minimum height=0.465cm] at (7, 4.25){} node[anchor=center, align=center, text width=0.077cm, inner sep=6pt] at (7, 4.25){$\overline{B}$};
	\node[shape=rectangle, minimum width=0.465cm, minimum height=0.465cm] at (2, 1.75){} node[anchor=center, align=center, text width=0.077cm, inner sep=6pt] at (2, 1.75){$\overline{A}$};
	\node[shape=rectangle, minimum width=0.465cm, minimum height=0.465cm] at (7, 1.75){} node[anchor=center, align=center, text width=0.077cm, inner sep=6pt] at (7, 1.75){$A$};
	\node[shape=rectangle, minimum width=0.465cm, minimum height=0.465cm] at (2, -1.25){} node[anchor=center, align=center, text width=0.077cm, inner sep=6pt] at (2, -1.25){$\overline{B}$};
	\node[shape=rectangle, minimum width=0.465cm, minimum height=0.465cm] at (7, -1.25){} node[anchor=center, align=center, text width=0.077cm, inner sep=6pt] at (7, -1.25){$B$};
	\node[pmos] at (0, 5.27){};
	\node[nmos] at (0, 3.73){};
	\draw (-0.98, 3.73) -| (-1, 4.5) |- (-0.98, 5.27);
	\draw (-1, 4.5) -- (-1.5, 4.5);
	\node[vcc] at (0, 6){};
	\node[ground] at (0, 2.96){};
	\node[pmos] at (0, 0.77){};
	\node[nmos] at (0, -0.77){};
	\draw (-0.98, -0.77) -| (-1, 0) |- (-0.98, 0.77);
	\draw (-1, 0) -- (-1.5, 0);
	\node[vcc] at (0, 1.5){};
	\node[ground] at (0, -1.54){};
	\node[shape=rectangle, minimum width=0.465cm, minimum height=0.465cm] at (-1.75, 4.5){} node[anchor=east, align=right, text width=0.077cm, inner sep=6pt] at (-1.5, 4.5){$A$};
	\node[shape=rectangle, minimum width=0.465cm, minimum height=0.465cm] at (-1.75, -0){} node[anchor=east, align=right, text width=0.077cm, inner sep=6pt] at (-1.5, -0){$B$};
	\node[jump crossing] at (-1, 4.25){};
	\node[jump crossing] at (0, 4.25){};
	\draw (-1.25, 4.5) -| (-1.25, 4.25) -- (-1.14, 4.25);
	\draw (-0.86, 4.25) -- (-0.14, 4.25);
	\draw (0, 4.5) -- (0.25, 4.5);
	\draw (0.14, 4.25) -- (0.25, 4.25);
	\node[jump crossing] at (-1, -0.25){};
	\node[jump crossing] at (-0, -0.25){};
	\draw (-1.25, -0) -| (-1.25, -0.25) -- (-1.14, -0.25);
	\draw (-0.86, -0.25) -- (-0.14, -0.25);
	\draw (-0, -0) -- (0.25, -0);
	\draw (0.14, -0.25) -- (0.25, -0.25);
	\draw (2.02, 5.77) -| (0.25, 5.75) -| (0.25, 4.5);
	\draw (0.25, 4.25) -| (0.5, 5) -- (2.86, 5);
	\draw (3.14, 5) -- (5.86, 5);
	\draw (6.14, 5) -- (7, 5) |- (6.98, 5.77);
	\node[jump crossing] at (3, 5){};
	\node[jump crossing] at (6, 5){};
	\draw (7, 5) -- (8, 5) |- (6.98, 1.25);
	\node[jump crossing] at (0.5, 4.5){};
	\draw (0.25, 4.5) -- (0.36, 4.5);
	\draw (2.04, 1.25) -- (0.75, 1.25) -- (0.75, 4.5) -- (0.64, 4.5);
	\draw (0, -0) -- (-1, -0);
	\node[jump crossing] at (3, -0){};
	\node[jump crossing] at (6, -0){};
	\draw (3.14, -0) -- (5.86, -0);
	\draw (6.14, -0) -- (8, -0) -| (8, -0.75) |- (6.98, -0.75);
	\node[jump crossing, rotate=90] at (1, 1.25){};
	\draw (1, 0.25) -- (1, 1.11);
	\draw (1, 1.39) -| (1, 3.75);
	\draw (2.02, 3.77) -| (1, 3.75);
	\draw (1.25, -0) -- (1.25, 0.5);
	\draw (1.25, 0.5) -- (7.75, 0.5) -| (7.75, 1);
	\node[jump crossing, rotate=90] at (7.75, 1.25){};
	\draw (7.75, 1) -| (7.75, 1.11);
	\draw (7.75, 1.39) -| (7.75, 3.75) |- (6.98, 3.77);
	\draw (1, 0.14) -| (1, 0.25);
	\draw (2.02, -0.75) |- (1.75, -0.75) -| (1.75, -0);
	\draw (1.75, -0) |- (0.25, -0);
	\draw (0.25, -0.25) -- (1, -0.25);
	\node[jump crossing, rotate=90] at (1, -0){};
	\draw (1, -0.25) -| (1, -0.14);
	\node[jump crossing] at (1.75, -0.25){};
	\draw (1, -0.25) -- (1.61, -0.25);
	\draw (1.89, -0.25) -- (2, -0.25) -| (2, -0) -- (2.86, -0);
\end{tikzpicture}
\end{center}

In generale, avremo bisogno di un inverter per variabile che compare in forma non standard nelle equazioni della PUN o della PDN.

Questo è esattamente il caso in cui la rete CMOS non è duale.

\subsection{Dimensionamento dei MOSFET}
Nella tecnologia CMOS il progettista della rete ha la possibilità di scegliere la dimensione del canale (cioè $W$ e $L$) di ogni transistor che usa.

Prendiamo l'esempio dell'inverter, e poniamo di avere i rapporti fra le dimensioni di canale $p$ ed $n$:
$$
\left( \frac{W}{L} \right)_P = p, \quad \left( \frac{W}{L} \right)_N = n
$$

I criteri che potremmo adottare sono:
\begin{enumerate}
\item \textbf{Area minima} di entrambi i canali: questo è semplice da implementare ma porta ad avere asimettrie nel comportamento dei due tipi di transistor;
\item \textbf{Simmetria}: in questo caso imponiamo l'equazione che avevamo trovato per la simmetria dell'inverter. L'equazione che ricordiamo volevamo rispettare per mantenere il bilanciamento fra porte PMOS e NMOS era:
$$
\frac{\mu_p}{\mu_n} = \frac{ \left( \frac{W}{L} \right)_N }{ \left( \frac{W}{L} \right)_P } = \frac{n}{p}
$$
Prendiamo quindi il rapporto fra le dimensioni dei canali come inverso del rapporto delle mobilità;
\item \textbf{Tempi di risposta}: fissato il \textit{rapporto} fra le dimensioni di canale, il parametro che resta da ottimizzare è il tempo di risposta $t_p$ della porta.
	\begin{itemize}
		\item Da una parte si può pensare di aumentare la dimensione di canale, in modo da aumentare la corrente condotta, e quindi la velocità di risposta;
		\item D'altro canto, aumentare la dimensione del canale aumenta la capacità equivalente, che ricordiamo era direttamente legata alla frequenza di operazione come:
			$$
				P_D = f C_{EQ} V_{DD}^2
			$$
	\end{itemize}
	La risposta ottima sarà quindi data da un processo di ottimizzazione a 2 obiettivi, che deve trovare un equilibrio fra potenza e capacità equivalente sulla porta.
\end{enumerate}

\par\smallskip

Riguardo alle \textbf{aree}, a questo punto, si avrà:
$$
A = W_N L_N + W_P L_P = n L_N^2 + p L_P^2
$$
Di cui si sceglie $L$ minima:
$$
L_N = L_P = L_{min}
$$
e si ricava il cosiddetto \textit{fattore di area} $(n + p)$:
$$
A = (n + p) L_{min}^2
$$

Una volta fissato l'invereter in una data famiglia logica, il dimensionamento delle reti PUN e PDN delle altre porte è fatto in maniera tale che nel caso peggiore queste abbiano lo stesso tempo di ritardo.

\par\smallskip

Possiamo fare delle considerazioni anche sulla \textit{potenza} considerando le \textbf{resistenze} associate ai transistori. Abbiamo in generale che la resistenza associata ad un MOSFET in conduzione fra drain e source (o viceversa nel PMOS) è inversamente proporzionale al rapporto $\frac{W}{L}$:
$$
R_{eq} = \frac{K}{\left(\frac{W}{L}\right)}
$$

A questo punto le resistenza equivalente di più porte sarà, se le porte sono in serie:
$$
R_{eq} = \sum_i R_i = \sum_i \frac{K}{\left(\frac{W}{L}\right)_i} = \frac{K}{\left(\frac{W}{L}\right)_{eq}} 
$$
e la $\frac{W}{L}$ equivalente sarà:
$$
\left( \frac{W}{L} \right)_{eq} = \left( \sum_i \frac{1}{ \left( \frac{W}{L} \right)_i }  \right)^{-1}
$$

\par\smallskip

\noindent
\textbf{\textsf{Dimensionamento di una porta NOR}}

\noindent
Iniziamo a vedere come si effettua il dimensionamento di una porta, prendendo l'esempio di una NOR a 4 ingressi.

\begin{center}
\begin{tikzpicture}[transform shape]
	% Paths, nodes and wires:
	\node[pmos] at (0, 0.77){};
	\node[pmos] at (0, 2.31){};
	\node[pmos] at (0, 3.85){};
	\node[pmos] at (0, 5.39){};
	\node[vcc] at (0, 6.16){};
	\node[nmos] at (-1, -0.77){};
	\node[nmos] at (-3, -0.77){};
	\node[nmos] at (1, -0.77){};
	\node[nmos] at (3, -0.77){};
	\draw (-3, -0) -- (3, -0);
	\node[ground] at (-3, -1.54){};
	\node[ground] at (-1, -1.54){};
	\node[ground] at (1, -1.54){};
	\node[ground] at (3, -1.54){};
	\draw[-stealth] (3, -0) -- (4, -0);
	\node[shape=rectangle, minimum width=0.465cm, minimum height=0.465cm] at (4.25, -0){} node[anchor=west, align=left, text width=0.077cm, inner sep=6pt] at (4, -0){$y$};
	\node[shape=rectangle, minimum width=0.465cm, minimum height=0.465cm] at (-1.23, 5.39){} node[anchor=east, align=right, text width=0.077cm, inner sep=6pt] at (-0.98, 5.39){$A$};
	\node[shape=rectangle, minimum width=0.465cm, minimum height=0.465cm] at (-4.23, -0.77){} node[anchor=east, align=right, text width=0.077cm, inner sep=6pt] at (-3.98, -0.77){$A$};
	\node[shape=rectangle, minimum width=0.465cm, minimum height=0.465cm] at (-1.23, 3.85){} node[anchor=east, align=right, text width=0.077cm, inner sep=6pt] at (-0.98, 3.85){$B$};
	\node[shape=rectangle, minimum width=0.465cm, minimum height=0.465cm] at (-2.23, -0.77){} node[anchor=east, align=right, text width=0.077cm, inner sep=6pt] at (-1.98, -0.77){$B$};
	\node[shape=rectangle, minimum width=0.465cm, minimum height=0.465cm] at (-1.23, 2.31){} node[anchor=east, align=right, text width=0.077cm, inner sep=6pt] at (-0.98, 2.31){$C$};
	\node[shape=rectangle, minimum width=0.465cm, minimum height=0.465cm] at (-0.23, -0.77){} node[anchor=east, align=right, text width=0.077cm, inner sep=6pt] at (0.02, -0.77){$C$};
	\node[shape=rectangle, minimum width=0.465cm, minimum height=0.465cm] at (-1.25, 0.75){} node[anchor=east, align=right, text width=0.077cm, inner sep=6pt] at (-1, 0.75){$D$};
	\node[shape=rectangle, minimum width=0.465cm, minimum height=0.465cm] at (1.77, -0.77){} node[anchor=east, align=right, text width=0.077cm, inner sep=6pt] at (2.02, -0.77){$D$};
\end{tikzpicture}
\end{center}

\begin{itemize}
\item Per quanto riguarda la \textbf{PUN}, abbiamo che la resistenza equivalente dei 4 transistor in serie deve essere uguale alla resistenza equivalente, per cui assumendoli uguali:
$$
\left( \frac{W}{L} \right)_{1, 2, 3, 4} = x
$$
si avrà:
$$
\frac{4}{x} = \frac{1}{\left( \frac{W}{L} \right)_{eq}} = \frac{1}{p} \implies x = 4p
$$
\item Per la \textbf{PDN}, invece, potremmo pensare di prendere la somma delle $\frac{W}{L}$. Dobbiamo però ricordare di dover prendere il \textit{caso peggiore}, cioè quello a resistenza massima: questo sarà quello in cui solo un transistor conduce. Assumiamo anche qui tutti i transistor uguali:
$$
\left( \frac{W}{L} \right)_{5, 6, 7, 8} = y
$$
da cui si avrà:
$$
\frac{1}{y} = \frac{1}{\left( \frac{W}{L} \right)_{eq}} = \frac{1}{n} \implies y = n
$$
\end{itemize}

L'area equivalente è data dlal fattore di area:
$$
A = (16 p + 4 n) L_{min}^2
$$

\par\smallskip

Il circuito duale alla NOR, cioè la NAND, sarà sostanzialmente simmetrico, ma con $x = p)$ delle porte PMOS e $y = 4n$ delle porte NMOS. Ora, visto che solitamente $\mu_n > \mu_p$ e quindi $n < p$, per la minimizzazione delle aree preferiamo usare la NAND, in quanto moltiplica per $16$ il termine più piccolo $n$.

In generale, il problema è che non vogliamo avere transistor in serie (in quanto per soddisfare i requisiti di tempi di risposta dobbiamo aumentare le dimensioni dei canali). Ora, se partiamo da avere transistor PMOS in serie, che già occupano più spazio (per la loro conducibilità minore), otteniamo dimensioni più grandi di quanto avremmo partendo da NMOS in serie (cioè il caso NOR contro NAND). Per questo motivo in tecnologia CMOS le NAND sono molto più diffuse delle NOR, e la sintesi dei circuiti si fa più volentieri in porte NAND. 

\subsection{Protezione dalle cariche elettrostatiche}
Se prendiamo il nostro inverter, abbiamo che a questo è associata una certa capacità equivalente. Se prendiamo l'equazione della capacità abbiamo:
$$
Q = CV \implies V = \frac{Q}{C}
$$
cioè la tensione fra le armature della capacità (per noi fra gate e substrato del MOSFET) è data dalla carica fratto la capacità. Visto che solitamente la capacità è molto piccola $\sim 10^{-15} \, \mathrm{F}$, ad una piccola carica corrisponde una grande tensione, che può danneggiare il MOSFET.

Per questo motivo dobbiamo predisporre dei circuiti di protezione, che anche in presenza del superamento della tensione in ingresso dell'inverter, evitano che questo si rompa. Questi saranno solitamente nella forma seguente:
\begin{center}
\begin{tikzpicture}[transform shape]
	% Paths, nodes and wires:
	\node[pmos] at (0.36, 0.77){};
	\node[nmos] at (0.36, -0.77){};
	\draw (-0.62, -0.77) -| (-0.64, 0) |- (-0.62, 0.77);
	\draw (-0.64, 0) -- (-1.14, 0);
	\node[vcc] at (0.36, 1.5){};
	\node[ground] at (0.36, -1.54){};
	\node[shape=rectangle, minimum width=0.465cm, minimum height=0.465cm] at (-4.25, -0){} node[anchor=east, align=right, text width=0.077cm, inner sep=6pt] at (-5, -0){$V_{in}$};
	\draw (0.36, 0) -- (0.61, 0);
	\draw (-2, -2) to[empty diode] (-2, -0);
	\draw (-2, -0) to[empty diode] (-2, 2);
	\draw (-2, -0) -| (-1.14, 0);
	\draw (-2, -0) to[american resistor, l={$R$}] (-4, -0);
	\node[vcc] at (-2, 2){};
	\node[ground] at (-2, -2){};
\end{tikzpicture}
\end{center}

La tensione (chiamiamola $V_K$) fra i diodi sarà compresa fra $0 \, \mathrm{V}$ e $V_{DD}$.
\begin{itemize}
	\item  Se $V_K$ prova a superare $V_{DD}$ di almeno $V_\gamma$, abbiamo che il diodo superiore entra in conduzione, e la $V_K$ viene fissata a $V_{DD} + V_\gamma$, cioè:
$$
V_K \geq V_{DD} + V_\gamma \implies V_K = V_{DD} + V_\gamma
$$
	\item Se d'altro canto $V_K$ prova a diminuire sotto lo $0$ di oltre $-V_\gamma$, abbiamo che il diodo inferiore entra in conduzione, e la $V_K$ viene fissata a $-V_\gamma$, cioé:
$$
V_K \leq - V_\gamma \implies V_K = - V_\gamma
$$
\end{itemize}

Notiamo che i circuiti di protezione vanno messi solo nei punti del circuito a contatto col mondo esterno. All'interno dei circuiti integrati, anzi, non è conveniente inserirli. Infatti, il diodo superiore è collegato al $V_{DD}$, per cui applica una corrente "a ritroso" pari a -$I_S$. Le correnti a ritroso, per quanto piccole, si accumulano su più porte in cascata, dandò quindi problemi di fan in.

\end{document}
